\chapter{第十一届大唐杯本科B组省赛第二场}
\fancyhead[L]{\color{H6}\kaishu\faIcon{atom}\;第十一届大唐杯本科B组省赛第二场}
\fancyhead[R]{\color{H6}\kaishu\rightmark\,}

\date{2024年4月20日}{}{\href{https://github.com/xubenshan/dtcup}{\textbf{才疏学浅的小熊}}}
{}
{https://www.xiaohongshu.com/user/profile/6192219d0000000021028647}{小红书}
{https://space.bilibili.com/692546609}{B站账号}



\section{单选题（共40小题，共120分）}



\begin{choice}{B}[]
    当前5G网络建设中，以下哪个频段分配给中国移动使用
    \begin{tasks}(2)
        \task 2715-2815Mhz
        \task 2515-2675Mhz
        \task 3400-3500Mhz
        \task 3500-3600Mhz
    \end{tasks}
\end{choice}


\begin{choice}{D}[]
    随机森林是决策树的集成，是一种什么方法
    \begin{tasks}(4)
        \task stacking
        \task Boosting
        \task pasting
        \task Bagging
    \end{tasks}
\end{choice}




\begin{choice}{D}[]
    根据香农公式，不能提高信道容量的方法
    \begin{tasks}(2)
        \task 使用多天线技术
        \task 采用更高效的调制方式
        \task 扩展信道带宽
        \task 降低信噪比
    \end{tasks}
\end{choice}



\begin{choice}{C}[]
    基站传输规划表中，30位掩码对应的子网掩码是
    \begin{tasks}(4)
        \task 255.255.255.240
        \task 255.255.255.248
        \task 255.255.255.252
        \task 255.255.255.0
    \end{tasks}
\end{choice}


\begin{choice}{B}[]
    产品开发评审主要有什么评审
    \begin{tasks}(2)
        \task 设计评审和测试评审
        \task 决策评审和技术评审
        \task 市场评审和生产评审
        \task 立项评审和开发评审
    \end{tasks}
\end{choice}

\begin{choice}{B}[]
    公司决策评审不包含
    \begin{tasks}(4)
        \task 发布决策评审
        \task 规格书评审
        \task 概念决策评审
        \task 计划决策评审
    \end{tasks}
\end{choice}




\begin{choice}{A}[]
    在5G通信系统中，信源编码的作用是提高通信的
    \begin{tasks}(4)
        \task 有效性
        \task 可靠性
        \task 安全性
        \task 经济性
    \end{tasks}
\end{choice}


\begin{choice}{A}[]
    关于摄像头、激光雷达、毫米波雷达三种传感器的比较，说法错误的是
    \begin{tasks}(1)
        \task 毫米波雷达的测距范围高于摄像头和激光雷达
        \task 激光雷达在目标检测精度上优于毫米波雷达和摄像头
        \task 摄像头的精度相对较低，视野范围最大
        \task 激光雷达的成本最高
    \end{tasks}
\end{choice}


\begin{choice}{B}[]
    5G\;NR系统中基于基站间的Xn链路切换，基站侧向核心网发的第一条信令为
    \begin{tasks}(2)
        \task Handover Request
        \task Handover Required
        \task SN Status Transfer
        \task Handover Request Acknowledge
    \end{tasks}
\end{choice}


\begin{choice}{C}[]
    移动通信每十年一代，6G技术是5G代际更新的一个新技术，预期其商用时间是
    \begin{tasks}(2)
        \task 2029年左右
        \task 2024年左右
        \task 2030年左右
        \task 2025年左右
    \end{tasks}
\end{choice}

\begin{choice}{C}[]
    在5G\;NR中AMF作为核心网主要的功能模块，以下属于AMF功能的描述为
    \begin{tasks}(2)
        \task PDU\;Session\;Control
        \task PDU\;Handing
        \task Idle\;State\;Mobility\;Handing
        \task Mobility\;Anchoring
    \end{tasks}
\end{choice}


\begin{choice}{B}[]
    下列哪个不是IPD的主流程
    \begin{tasks}(2)
        \task 市场规划流程
        \task 物料申请流程
        \task 技术开发流程
        \task 产品开发流程
    \end{tasks}
\end{choice}


\begin{choice}{D}[]
    信息增益在决策树算法中是用来选择特征的指标，生成树时先选择信息增益（\quad）的特征
    \begin{tasks}(4)
        \task 小
        \task 以上都不是
        \task 平均值
        \task 大
    \end{tasks}
\end{choice}


\begin{choice}{D}[]
    以下不属于决策树求解算法的选项是
    \begin{tasks}(4)
        \task ID3
        \task C4.5
        \task CART算法
        \task 梯度下降法
    \end{tasks}
\end{choice}

\begin{choice}{D}[]
    3GPP未来将向6G逐步演进，3GPP的NTN标准是从哪一年开始启动，一直朝着将卫星纳入3GPP技术规范的目标在前进
    \begin{tasks}(4)
        \task R16
        \task R18
        \task R17
        \task R15
    \end{tasks}
\end{choice}

\begin{choice}{A}[]
    若SCS为30kHz，则一个RB的宽度是
    \begin{tasks}(4)
        \task 360
        \task 1440
        \task 180
        \task 720
    \end{tasks}
\end{choice}



\begin{choice}{D}[]
    5G车联网技术在智慧交通系统中的作用不包括
    \begin{tasks}(4)
        \task 车辆定位与导航
        \task 紧急情况下的通信
        \task 实时路况监控
        \task 车辆娱乐技术
    \end{tasks}
\end{choice}

\begin{choice}{A}[]
    3GPP协议针对5G第一个商用版本为
    \begin{tasks}(4)
        \task R15
        \task R14
        \task R16
        \task R17
    \end{tasks}
\end{choice}


\begin{choice}{C}[]
    在5G基站安装中，DCPD与电源柜之间的电源线规格应为
    \begin{tasks}(4)
        \task 10方
        \task 16方
        \task 25方
        \task 8方
    \end{tasks}
\end{choice}


\begin{choice}{C}[]
    RRC-INACTIVE状态被称为挂起状态，如果恢复业务，则会出现哪个信令
    \begin{tasks}(2)
        \task RRC connection request
        \task RRC connection reconfiguration
        \task RRC resume request
        \task RRC connection reject
    \end{tasks}
\end{choice}



\begin{choice}{D}[]
    华为Mate60pro的卫星通话功能采用中国自主研制的卫星移动通信系统由谁运营
    \begin{tasks}(2)
        \task 中国移动
        \task 中国卫通
        \task 中国联通
        \task 中国电信
    \end{tasks}
\end{choice}




\begin{choice}{A}[]
    5G网络在理想前传条件下，DU可选择的部署方式为
    \begin{tasks}(4)
        \task 集中部署
        \task 一体化
        \task 远端部署
        \task 随意部署
    \end{tasks}
\end{choice}

\begin{choice}{D}[]
    预计6G通信将支持五个应用场景，其中安全可靠低时延通信场景可简称为
    \begin{tasks}(4)
        \task eMBB-Plus
        \task 3D-InteCom
        \task BigCom
        \task SURLLC
    \end{tasks}
\end{choice}

\begin{choice}{D}[]
    DC供电系统的电压波动范围为
    \begin{tasks}(2)
        \task -40V至-53V
        \task -40V至-55V
        \task -40V至-50V
        \task -40V至-57V
    \end{tasks}
\end{choice}


\begin{choice}{B}[]
    6G网络将在智享生活，智赋生产等方面催生全新的应用场景，这些场景将需要很多全新的网络功能和能力，其中不包括
    \begin{tasks}(2)
        \task 数字孪生
        \task 毫秒级的时延体验
        \task 太比特级的峰值速率
        \task 超过1000km/h的移动速度
    \end{tasks}
\end{choice}


\begin{choice}{C}[]
    初始随机接入过程中，表示竞争解决的消息为
    \begin{tasks}(4)
        \task MSG5
        \task MSG2
        \task MSG4
        \task MSG3
    \end{tasks}
\end{choice}



\begin{choice}{D}[]
    无线对讲通话设备主要用于提供实时语音服务，其主要作用不包括
    \begin{tasks}(4)
        \task 应急通信
        \task 团队协作
        \task 便捷通信
        \task 抗干扰通信
    \end{tasks}
\end{choice}


\begin{choice}{A}[]
    在数据存储阶段，主要困难是由于数据的(\quad)特性造成数据的分类管理和防护难度大
    \begin{tasks}(4)
        \task 多样性
        \task 简单化
        \task 复杂化
        \task 格式化
    \end{tasks}
\end{choice}


\begin{choice}{D}[]
    5G中PBCH在频域有(\quad)个SubCarrier
    \begin{tasks}(4)
        \task 144
        \task 48
        \task 127
        \task 240
    \end{tasks}
\end{choice}


\begin{choice}{C}[]
    NR系统中，UE可通过哪个系统消息获取小区重选公共信息
    \begin{tasks}(4)
        \task SIB3
        \task SIB4
        \task SIB2
        \task SIB1
    \end{tasks}
\end{choice}


\begin{choice}{D}[]
    毫米波雷达是一种重要的ADAS传感器，可以提供多种高精度的路面空间信息，其中，不包括
    \begin{tasks}(4)
        \task 方位角
        \task 相对速度
        \task 距离
        \task 目标分类
    \end{tasks}
\end{choice}


\begin{choice}{C}[]
    以下哪个方法适用于对生产周期短、生产数量大的产品进行成本降低、成本控制的管理
    \begin{tasks}(2)
        \task 作业变动成本法
        \task 完全成本法
        \task 变动成本法
        \task 作业成本法
    \end{tasks}
\end{choice}



\begin{choice}{D}[]
    在5G NR中定义一个子帧时间长度是
    \begin{tasks}(4)
        \task 5ms
        \task 10ms
        \task 0.5ms
        \task 1ms
    \end{tasks}
\end{choice}


\begin{choice}{D}[]
    在5G NR有多少个唯一的物理层小区ID
    \begin{tasks}(4)
        \task 504
        \task 256
        \task 512
        \task 1008
    \end{tasks}
\end{choice}


\begin{choice}{B}[]
    ModBus标准，主要解决一根双绞线实现与（\quad）个设备间的通信，也是当前最受欢迎的工业现场通信协议的一种
    \begin{tasks}(4)
        \task 无数
        \task 多
        \task 一
        \task 无限
    \end{tasks}
\end{choice}


\begin{choice}{D}[]
    Massive MIMO是5G关键技术之一，是解决什么需求的有效技术
    \begin{tasks}(4)
        \task 低时延
        \task 高清画质
        \task 智能应用
        \task 大容量
    \end{tasks}
\end{choice}


\begin{choice}{D}[]
    5G NR协议栈中，以下选项中属于5G空口控制面层二协议栈
    \begin{tasks}(4)
        \task SCTP
        \task UDP
        \task SDAP
        \task PDCP
    \end{tasks}
\end{choice}


\begin{choice}{D}[]
    多个GPS天线之间的间距不能过近，否则天线之间可能会产生反射干扰，两个GPS天线间距大于
    \begin{tasks}(4)
        \task 0.5
        \task 4
        \task 1
        \task 2
    \end{tasks}
\end{choice}


\begin{choice}{D}[]
    高压线附近建设站点时，通信机房应保持（\quad）m以上的距离
    \begin{tasks}(4)
        \task 10
        \task 5
        \task 20
        \task 15
    \end{tasks}
\end{choice}


\begin{choice}{B}[]
    已知项目的计息周期为月，月利率为千分之8，则年实际利率为
    \begin{tasks}(4)
        \task 0.08
        \task 0.1003
        \task 千分之9.6
        \task 0.096
    \end{tasks}
\end{choice}



\section{多选题（共30小题，共120分）}

\begin{choice}{\;BC\;}[]
    移动通信系统中，影响无线电波传播路径损耗的基本传播机制为
    \begin{tasks}(4)
        \task 直射
        \task 反射
        \task 绕射
        \task 折射
    \end{tasks}
\end{choice}

\begin{choice}{\;ABC\;}[]
    5GD2D通信优势包括
    \begin{tasks}(2)
        \task 短距离通信可频谱资源复用
        \task 终端近距离通信，高速率低时延低功耗
        \task 拓展网络覆盖范围
        \task 无线M2M功能
    \end{tasks}
\end{choice}

\begin{choice}{\;ABCD\;}[]
    根据通信系统的模型特点，以下属于通信系统一般模型中基本组成部分的有
    \begin{tasks}(4)
        \task 信道
        \task 信宿
        \task 发送和接收设备
        \task 信源
    \end{tasks}
\end{choice}

\begin{choice}{\;ABD\;}[]
    下列选项中属于5G参考信号的是
    \begin{tasks}(4)
        \task SRS
        \task DMRS
        \task CRS
        \task PTRS
    \end{tasks}
\end{choice}


\begin{choice}{\;ABD\;}[]
    中国提出的非正交多址技术有哪些
    \begin{tasks}(4)
        \task SCMA
        \task   PDMA
        \task NOMA
        \task MUSA
    \end{tasks}
\end{choice}

\begin{choice}{\;ABD\;}[]
    专利应对机制的策略包括
    \begin{tasks}(2)
        \task 专利避让
        \task 交纳专利使用费
        \task 以上均不是
        \task 交叉许可
    \end{tasks}
\end{choice}

\begin{choice}{\;BD\;}[]
    以下关于6G网络性能指标初步预测情况正确的是
    \begin{tasks}(2)
        \task 定位精度10m
        \task UP时延小于0.5ms
        \task 用户体验速率100Mbps
        \task 峰值数据率1Tbps
    \end{tasks}
\end{choice}


\begin{choice}{\;ABCD\;}[]
    车联网中，环境感知主要包括
    \begin{tasks}(1)
        \task 车辆本身状态感知，包括行驶速度，行驶方向，行驶状态，车辆位置等
        \task 交通信号感知，自动识别交叉路口的信号灯，如何高效通过交叉路口等
        \task 道路感知、包括道路类型检测，道路标线识别，道路状况判断，是否偏离行驶轨迹等
        \task 行人感知，主要判断车辆行驶前方是否有行人，包括白天行人识别，夜晚行人识别，被障碍物遮挡的行人识别
    \end{tasks}
\end{choice}


\begin{choice}{\;ABCD\;}[]
    在5G中CSS（公共搜索空间）有多种类型，不同的类型对应不同的消息，以下属于CSS分类的为
    \begin{tasks}(2)
        \task Type3 PDCCH CSS
        \task Type0A PDCCH CSS
        \task Type2 PDCCH CSS
        \task Type1 PDCCH CSS
    \end{tasks}
\end{choice}


\begin{choice}{\;ABCD\;}[]
    在5G基站日常维护过程中，导致AC校准问题的原因可能是
    \begin{tasks}(2)
        \task 驻波比问题
        \task 输出功率问题
        \task GPS时钟源异常问题
        \task 光链路异常问题
    \end{tasks}
\end{choice}


\begin{choice}{\;ABCD\;}[]
    根据香农公式，提高信道容量主要有以下方法

    \begin{tasks}(2)
        \task 降低信号干扰
        \task 提升信噪比
        \task 扩展信道带宽
        \task 采用更高效的调制方式
    \end{tasks}
\end{choice}

\begin{choice}{\;ABCD\;}[]
    下面属于集成学习方法的选项是
    \begin{tasks}(4)
        \task Boosting
        \task stacking
        \task Bagging
        \task Voting
    \end{tasks}
\end{choice}

\begin{choice}{\;AB\;}[]
    在5G NR中，PDCCH时频域位置的确定与哪些因素有关
    \begin{tasks}(4)
        \task Search Space
        \task CORESET
        \task UCI
        \task DCI
    \end{tasks}
\end{choice}

\begin{choice}{\;ABC\;}[]
    在5G网络中，AMF通过N2接口连接基站，以下属于N2接口协议栈的为
    \begin{tasks}(4)
        \task GTP-U
        \task NGAP
        \task SCTP
        \task SDAP
    \end{tasks}
\end{choice}

\begin{choice}{\;ABD\;}[]
    5G NR系统中毫米波通信的优势有
    \begin{tasks}(2)
        \task 可用频带宽，可提供几十GHz带宽
        \task 方向性好、受干扰影响小
        \task  绕射能力强、路径损耗小
        \task 波束集中，提高能效
    \end{tasks}
\end{choice}

\begin{choice}{\;ABCD\;}[]
    经济学中资源的基本形态包括
    \begin{tasks}(2)
        \task 人（劳动）
        \task 土地和自然资源
        \task 资本（机器设备、厂房等生产工具）
        \task 企业家才能
    \end{tasks}
\end{choice}

\begin{choice}{\;AC\;}[]
    以下说法正确的是
    \begin{tasks}(1)
        \task 当新兴市场的战略机会显现，所有人都能看到时，就代表着错失良机
        \task 商业机会的把握其实靠运气，并无规律可循
        \task 即使技术不先进，商业模式选择正确也可能盈利
        \task 只要技术足够创新，足够先进，就一定会盈利
    \end{tasks}
\end{choice}

\begin{choice}{\;ABD\;}[]
    我国的固定信令系统N0.7信令网采用的级别的分类包括
    \begin{tasks}(1)
        \task 高级信令转接点HSTP、汇接LSTP和所属的SP
        \task 信令点SP
        \task 信令转接点STP
        \task 汇接转接点LSTP、汇接所属的SP
    \end{tasks}
\end{choice}

\begin{choice}{\;ABD\;}[]
    移动通信系统中，关于扩频码和扰码说法不正确的是
    \begin{tasks}(1)
        \task 扩频码不需要改变带宽，扰码不改变信号带宽
        \task 扩频码不需要改变带宽，扰码需要改变信号带宽
        \task 扩频码需要改变带宽，扰码不改变信号带宽
        \task 扩频码需要改变带宽，扰码需要改变信号带宽
    \end{tasks}
\end{choice}

\begin{choice}{\;ACD\;}[]
    以下属于决策树求解算法的选项是
    \begin{tasks}(4)
        \task CART算法
        \task 梯度下降法
        \task C4.5
        \task ID3
    \end{tasks}
\end{choice}

\begin{choice}{\;AB\;}[]
    关于工业互联网的发展和应用，以下说法正确的是
    \begin{tasks}(1)
        \task 随着5G技术的发展，工业互联网的传输速度和可靠性将得到显著提升
        \task 工业互联网的推广有助于实现定制化生产和服务，满足消费者的个性化需求
        \task 工业互联网的实施不需要考虑现有设备的兼容性，因为新技术可以完全替代旧设备
        \task 工业互联网的应用仅限于大型企业，中小型企业无法从中受益
    \end{tasks}
\end{choice}

\begin{choice}{\;ABCD\;}[]
    5GFR2最大传输带宽配置
    \begin{tasks}(4)
        \task 100Mhz
        \task 50Mhz
        \task 400Mhz
        \task 200Mhz
    \end{tasks}
\end{choice}

\begin{choice}{\;AD\;}[]
    在5GPDU会话建立过程中，PDU\;Session\;Resource\;setup\;request包含哪些信息
    \begin{tasks}(2)
        \task PDU session Aggregate Maximum Bitrate
        \task PDU-UE-NGAP-ID
        \task 终端业务ip地址
        \task Qos\;id
    \end{tasks}
\end{choice}


\begin{choice}{\;ABCD\;}[]
    6G技术的发展基于信息普惠、绿色能耗、虚拟融合三大发展理念，以下属于6G前沿技术分类的有
    \begin{tasks}(2)
        \task 天线人工智能
        \task 智能超表面
        \task 太赫兹和可见光通信
        \task 全息无线电
    \end{tasks}
\end{choice}


\begin{choice}{\;BCD\;}[]
    大唐5G基站设备维护是，下列那些原因可能导致AAU发送功率异常
    \begin{tasks}(4)
        \task UE类型
        \task 滤波器故障
        \task 通道故障
        \task AC系数
    \end{tasks}
\end{choice}


\begin{choice}{\;ABCD\;}[]
    工业互联网平台应用层的作用是
    \begin{tasks}(2)
        \task 提供工业微服务
        \task 工业数据管理
        \task 业务应用
        \task APP开发
    \end{tasks}
\end{choice}


\begin{choice}{\;ACD\;}[]
    智能网联汽车环境感知传感器属于信息收集单元的重要组成部分，其信息收集能力主要有什么决定
    \begin{tasks}(2)
        \task 传感器感知精度
        \task 车辆所在位置
        \task 传感器的数量
        \task 传感器感知范围
    \end{tasks}
\end{choice}


\begin{choice}{\;BCD\;}[]
    经营成本是指总成本费用扣除什么以后的全部费用
    \begin{tasks}(4)
        \task 销售税金及附加
        \task 折旧费
        \task 摊销费
        \task 利息支出
    \end{tasks}
\end{choice}


\begin{choice}{\;ABCD\;}[]
    从5G到6G，超大规模天线技术持续升级，其中什么是其未来重大研究方向
    \begin{tasks}(2)
        \task 精准定位
        \task 超维度无边界链接
        \task 智能超表面RIS
        \task 电磁波角动能
    \end{tasks}
\end{choice}


\begin{choice}{\;ABD\;}[]
    5G系统中，SS/PBCH\;block的子载波间隔$\mu $值可以配置为
    \begin{tasks}(4)
        \task 3
        \task 1
        \task 2
        \task 0
    \end{tasks}
\end{choice}


\section{判断题（共30小题，共60分）}

\begin{choice}{\;正确\;}[]
    在数值上，载噪比要比载干比大一些。
\end{choice}


\begin{choice}{\;错误\;}[]
    5G的基站功能重构为CU和DU两个功能实体，CU主要处理物理层功能和实时性需求的层2功能。
\end{choice}

\begin{choice}{\;正确\;}[]
    产品研发时不仅要考虑研发成本，还要考虑全周期、全过程的成本。
\end{choice}

\begin{choice}{\;正确\;}[]
    GBDT不像AdaBoost那样在每个迭代中调整实例权重，而是让新的预测器针对前一个预测器的残差进行拟合。
\end{choice}

\begin{choice}{\;错误\;}[]
    车联网环境中，两车相距较远时，相互之间可以通过Pc5接口通信。
\end{choice}

\begin{choice}{\;错误\;}[]
    固定电话网的干线传输设备中，以缆线传输设备为主，其他传输设备为辅。

\end{choice}


\begin{choice}{\;正确\;}[]
    6G网络是一种概念性的无线网络移动通信技术，也叫做第六代移动通信技术。

\end{choice}

\begin{choice}{\;正确\;}[]

    5G NR系统中NG、Xn、F1接口信令连接都基于GTP-U协议。
\end{choice}


\begin{choice}{\;正确\;}[]
    26Ghz为厘米波，可用频带宽，可提供几十GHz带宽，提供高速用户体验，更适合微站补热场景。

\end{choice}

\begin{choice}{\;正确\;}[]
    标准是工业互联网的法则，工业互联网标准化工作是实现工业互联网的重要技术基础。

\end{choice}


\begin{choice}{\;正确\;}[]
    5G NR系统中随机接入类型分为竞争和非竞争两种。

\end{choice}


\begin{choice}{\;正确\;}[]
    华为Mate60pro的卫星通话功能采用中国自主研制的卫星移动通信系统由中国电信运营。

\end{choice}


\begin{choice}{\;错误\;}[]
    在塔上安装GPS天线，应使GPS天线位于铁塔最面向南方一角并保持与铁塔0.5米间距。

\end{choice}


\begin{choice}{\;正确\;}[]
    通信系统的作用就是将信息从信源发送到一个或多个目的地。

\end{choice}


\begin{choice}{\;错误\;}[]
    移动通信系统中，多径效应和多普勒频移可以引起慢衰落。

\end{choice}

\begin{choice}{\;正确\;}[]
    位于山坡上的城镇覆盖，基站宜设置在山顶。

\end{choice}


\begin{choice}{\;正确\;}[]
    在交叉路口中，车联网汽车可以通过V2V方式获取路口信息，降低碰撞风险。

\end{choice}


\begin{choice}{\;错误\;}[]
    产品变动成本是与产品紧密联系的，产品总的变动成本与产品数量是线性关系。

\end{choice}


\begin{choice}{\;正确\;}[]
    5G组网方式可分为独立部署和非独立部署两种。

\end{choice}


\begin{choice}{\;正确\;}[]
    在位置满足要求的情况下，GPS到接收机的馈线可以长一些减少成本支出。

\end{choice}


\begin{choice}{\;正确\;}[]
    5G NR中SSB的频点是基于其占用的20个PRB中的第11个PRB的子载波0中心频率。

\end{choice}


\begin{choice}{\;正确\;}[]
    工业互联网数据安全，指工厂内部重要的生产管理数据，生产操作数据以及工厂外部数据等各类数据的安全。

\end{choice}


\begin{choice}{\;错误\;}[]
    5G NR系统中，UE在RRC IDLE态RRC INACTIVE态的寻呼标识为P-RNTI。

\end{choice}


\begin{choice}{\;错误\;}[]
    熵代表随机变量的确定性。

\end{choice}

\begin{choice}{\;正确\;}[]
    在3GPP\;R17中，5G技术在支持人工智能和机器学习应用方面的重要性和潜力得到了重点强调。

\end{choice}

\begin{choice}{\;错误\;}[]
    根据香农公式，在要求信道容量一定的情况下，信道带宽越大，对信噪比的要求则越高。

\end{choice}

\begin{choice}{\;正确\;}[]
    在NR中为了标识射频信道、SSB或者其他资源的频域位置而定义的基本单位为全局栅格。

\end{choice}

\begin{choice}{\;错误\;}[]
    完全成本法是管理会计主要采用的成本核算方法。

\end{choice}

\begin{choice}{\;错误\;}[]
    5G NR中信令承载DRB被定义为用于传递RRC和NAS消息的无线承载。

\end{choice}

\begin{choice}{\;错误\;}[]
    AAU和BBU间的数据传输称为回传。

\end{choice}





